% Independent finite compatibility derivation. Prefix MRR.
% Requires amsmath, amssymb and hyperref. All original metrics retained.
\section{Exact monodromy compatibility of the original prime receiver}
\label{sec:monodromy-receiver-review}

\paragraph{Source convention and original input.}
Uwe Jannsen, \emph{Weights in Arithmetic Geometry},
\href{https://arxiv.org/abs/1003.0927v1}{arXiv:1003.0927v1},
original author file \nolinkurl{main.tex}, lines 325--365, was read
for this calculation. Lines 335--340 state
$\mathcal N F=q_{\rm res}F\mathcal N$ for geometric Frobenius
and nilpotent monodromy, with $q_{\rm res}$ the residue-field
cardinality. Lines 341--352 specify the associated monodromy
filtration and twists. The retained source has SHA-256
\par\noindent
\nolinkurl{deadee9fddd84a96a5f8a5ad282f929ad837b1f0a67bbb04907980681b12ab0a}.
\par\noindent
We test the stated geometric convention with residue field
$\mathbb F_p$, so its equation becomes
$F_p\mathcal N F_p^{-1}=p^{-1}\mathcal N$.
The programme's symbol $q=(k+1)^2$ remains its vector-space
dimension and is not the $q_{\rm res}$ in Jannsen's equation.
This is a finite compatibility calculation, not an assertion
that the programme operators already have an arithmetic realization.

The precise programme input is ST1--ST3 of the retained
\nolinkurl{SUPPORT_TRANSPORT_PUBLIC.md} and TDP1--TDP4, TDP11
of \nolinkurl{TAU_DELiGNE_PRIME_MEMORY.tex}.
ST3 states the exact source-metric decomposition below and
locates its proof in
\href{https://github.com/KokunoYumeto/zeta-function-research-reader/blob/5c69161ca70ee187f41df0bfe786e8b6eebed422/workbenches/splitzero-tandem/continuations/20260922-current-effective-cutoff/EVOLUTION_PROOF.md#L18}{CE2--CE8
of the pinned original effective-cutoff proof}.
The present derivation read ST1--ST3 and their displayed proof
locator; it does not claim a new derivation of those source
identities from an unspecified metric.
Keep every original parameter and cutoff:
\[
\begin{gathered}
k\ge17,\quad k\equiv1\pmod4,\quad q=(k+1)^2,
\quad\delta,\gamma>0,\quad q-1\le N\le2q,\\
\lambda_{ab}=(2b-k)\gamma-i(2a-k)\delta,\quad
Q_k(y)=\prod_{a,b=0}^k(y-\lambda_{ab}),\quad
E=\mathbb C[y]/(Q_k),\quad M[v]=[yv],\\
F_p=e^{i(\log p)M},\quad\widetilde F_p=p^{k/2}F_p,
\quad \langle v,w\rangle_N=v^*G_Nw,\\
M=C+R,\quad C=C^\dagger,\quad R=\epsilon fe^\dagger,
\quad e^\dagger f=0,\quad\|e\|_N=\|f\|_N=1,\quad\epsilon>0.
\end{gathered}
\tag{MRR1}
\]
Here the actual ST3 vectors are
$e=b_N/\sqrt{E_N}$, $f=b_{N+1}/\sqrt{F_N}$ and
$\epsilon=\sqrt{E_NF_N}/\omega_N$. These are the retained
source definitions, not substitutions made in this derivation.
The four cutoffs $q-1,q,2q-1,2q$ remain distinct.

\subsection{Two primes determine the complete all-prime space}

Let $\ell_{ab}$ be the original Lagrange cardinal polynomials.
Their distinct roots give a basis of $E$ by evaluation, and
$M\ell_{ab}=\lambda_{ab}\ell_{ab}$.
Let $\ell^{ab}$ denote the linear evaluation covector at
$\lambda_{ab}$, so $\ell^{ab}(\ell_{cd})=
\mathbf1_{(a,b)=(c,d)}$. It is not silently replaced by a
Hermitian adjoint. Put
$U_{ab,cd}=\ell_{ab}\ell^{cd}$ and
$P_{ab}=U_{ab,ab}$.
These are the exact matrix units and spectral projectors in
the original coordinates. In particular
\[
F_pU_{ab,cd}F_p^{-1}
=p^{2(a-c)\delta}
 e^{2i(b-d)\gamma\log p}U_{ab,cd},\qquad
[M,U_{ab,cd}]=(\lambda_{ab}-\lambda_{cd})U_{ab,cd}.
\tag{MRR2}
\]
Applying either side to every original cardinal vector proves
these identities, including their phases.

The complete compatibility space is
\[
\begin{split}
\mathscr N
&=\{T:F_2TF_2^{-1}=2^{-1}T,
            \ F_3TF_3^{-1}=3^{-1}T\}\\
&=\{T:F_pTF_p^{-1}=p^{-1}T\text{ for every prime }p\}\\
&=\ker(\operatorname{ad}_M-iI).
\end{split}
\tag{MRR3}
\]
To prove equality, a nonzero matrix-unit coefficient must
satisfy both scalar equations in MRR2. Their moduli force
$2(a-c)\delta=-1$. Their phases force
$2(b-d)\gamma\log2=2\pi n_2$ and
$2(b-d)\gamma\log3=2\pi n_3$, with integers $n_2,n_3$.
If $(b-d)\gamma\ne0$, both integers are nonzero, so
$\log2/\log3=n_2/n_3$ would be rational.
This is impossible: a rational equality gives, after choosing
positive numerator and denominator, $2^v=3^u$, contradicted
by the prime divisors of these integers.
Thus $b=d$ because $\gamma>0$, and the modulus equality gives
$\lambda_{ab}-\lambda_{cd}=i$.
Conversely this eigenvalue difference makes the scalar in
MRR2 equal $e^{-\log p}=p^{-1}$ for every prime, and also
gives the commutator equation. Independence of all matrix
units proves the equalities for every operator.

Writing $m=c-a$, the space is therefore exactly
\[
\boxed{
\begin{cases}
\mathscr N=0,&\delta\ne1/(2m)\text{ for all }1\le m\le k,\\[2pt]
\mathscr N=
 \displaystyle\bigoplus_{a=0}^{k-m}\bigoplus_{b=0}^k
          \mathbb C U_{ab,a+m,b},
 &\delta=1/(2m),\quad1\le m\le k.
\end{cases}}
\tag{MRR4}
\]
Indeed $a,c$ are integers in $[0,k]$ and $\delta>0$, so the
equation $2(a-c)\delta=-1$ has precisely that range of positive
$m$. For a fixed $\delta$ at most one $m$ occurs.
At resonance its dimension is $(k+1-m)(k+1)$, and otherwise
it is zero.
Multiplication of $F_p$ by its retained scalar $p^{k/2}$ does
not change conjugation, so exactly the same space results
from the original uncentred $\widetilde F_p$.

Every operator in this space is nilpotent, with its whole
chain law specified. At resonance write
$T=\sum_{a=0}^{k-m}\sum_b n_{ab}U_{ab,a+m,b}$.
Then
\[
T^s\ell_{ab}=
\begin{cases}
\left(\displaystyle\prod_{j=1}^s n_{a-jm,b}\right)
                         \ell_{a-sm,b},&a\ge sm,\\
0,&a<sm,
\end{cases}
\quad T^{\lfloor k/m\rfloor+1}=0.
\tag{MRR5}
\]
The first step is the definition, and induction proves the
product for every $s$. It lowers $a$ by $m$ and keeps $b$
fixed, so the last equality follows from $a\le k$.
The nilpotence bound is sharp: taking all $n_{ab}=1$ keeps
the longest chain from $a=k$ nonzero through
$\lfloor k/m\rfloor$ steps.

\subsection{The exact spectral and metric projections}

At resonance the coordinate projection is
\[
\mathsf S(T)=\sum_{a=0}^{k-m}\sum_{b=0}^k
                 P_{ab}T P_{a+m,b};
\qquad \mathsf S=0\text{ off resonance}.
\tag{MRR6}
\]
The matrix-unit relations prove $\mathsf S^2=\mathsf S$ and
$\operatorname{im}\mathsf S=\mathscr N$.
It commutes with $\operatorname{ad}_M$ and all original prime
conjugations because it selects their simultaneous coordinate
eigenvectors. It is not asserted to be orthogonal in $G_N$.

The actual metric projection is also explicit.
Let $V$ be the invertible matrix whose columns are all
original $\ell_{ab}$ in their retained ordering, and put
$H=V^*G_NV$.
On $\operatorname{End}(E)$ use the exact Hilbert--Schmidt
inner product
$\langle A,B\rangle_{\rm HS,N}=\operatorname{tr}(A^\dagger B)$,
$A^\dagger=G_N^{-1}A^*G_N$.
In the cardinal coordinates this is
$\operatorname{tr}(H^{-1}A^*HB)$.
For matrix units labelled by pairs $(i,j)$, $(u,v)$, the
complete Gram entries are
\[
\langle U_{ij},U_{uv}\rangle_{\rm HS,N}
                          =H_{iu}(H^{-1})_{vj}.
\tag{MRR7}
\]
Indeed $U_{ij}^\dagger=H^{-1}U_{ji}H$ in those coordinates;
the product $U_{ji}HU_{uv}$ is $H_{iu}U_{jv}$, whose trace
after multiplication by $H^{-1}$ is the displayed entry.
Thus every off-diagonal metric contribution is retained.

Let $\{K_\alpha\}$ be exactly the compatible matrix units
in MRR4 and define
$\Gamma_{\alpha\beta}=\langle K_\alpha,K_\beta\rangle_{\rm HS,N}$,
$b_\alpha(T)=\langle K_\alpha,T\rangle_{\rm HS,N}$.
Their independence and the positive metric make $\Gamma$
positive definite. The exact orthogonal projection and
distance are
\[
\boxed{
\begin{aligned}
\mathsf P_N T
 &=\sum_{\alpha,\beta}K_\alpha
                   (\Gamma^{-1})_{\alpha\beta}b_\beta(T),\\
\operatorname{dist}_{\rm HS,N}(T,\mathscr N)^2
 &=\|T\|_{\rm HS,N}^2-b(T)^*\Gamma^{-1}b(T).
\end{aligned}}
\tag{MRR8}
\]
The coefficient vector solves precisely
$\langle K_\alpha,T-\mathsf P_NT\rangle=0$ for every $\alpha$.
This proves orthogonality and the norm formula by the
Pythagorean identity. Off resonance both projections vanish,
and the distance is $\|T\|_{\rm HS,N}$.

\subsection{The complete two-prime defect and its inverse receiver}

For $\mathcal E=\operatorname{End}(E)$ define
\[
\begin{gathered}
\mathcal D_0(T)=[M,T]-iT,\qquad
\mathcal D_p(T)=F_pTF_p^{-1}-p^{-1}T,\qquad
\mathcal D_{2,3}(T)=(\mathcal D_2(T),\mathcal D_3(T)),\\
\mathcal E/\mathscr N\xrightarrow{\ \simeq\ }
             \operatorname{im}\mathcal D_{2,3},
\qquad [T]\longmapsto\mathcal D_{2,3}(T).
\end{gathered}
\tag{MRR9}
\]
MRR3 proves that the last map is well-defined and injective;
surjectivity onto its specified image is its definition.
The analogous map $[T]\mapsto\mathcal D_0(T)$ is likewise
an isomorphism onto its image.
These are the exact defect objects even when a nonzero
compatible space survives.

Write
$\mu_{ij}=\lambda_i-\lambda_j-i$ and
$\xi_{p,ij}=\omega_{p,i}/\omega_{p,j}-p^{-1}$, where
$\omega_{p,ab}=p^{(2a-k)\delta}e^{i(2b-k)\gamma\log p}$.
Both maps are diagonal on the same matrix units, with
coefficients $\mu_{ij}$ and $\xi_{p,ij}$ respectively.
Their full original-metric energy is therefore
\[
\begin{split}
\|\mathcal D_{2,3}(T)\|_{\rm HS,N\oplus HS,N}^2
=\sum_{ij,uv}\overline{t_{ij}}t_{uv}
H_{iu}(H^{-1})_{vj}
\bigl(\overline{\xi_{2,ij}}\xi_{2,uv}
     +\overline{\xi_{3,ij}}\xi_{3,uv}\bigr),
\quad T=\sum_{ij}t_{ij}U_{ij}.
\end{split}
\tag{MRR10}
\]
This is expansion of each squared norm using MRR7, with no
assumption that the matrix units are orthogonal.

There is an entire forward relation and an exact finite
reverse relation:
\[
\begin{gathered}
\mathcal D_p=g_p(\mathcal D_0)\mathcal D_0,\qquad
g_p(z)=p^{-1}\frac{e^{i(\log p)z}-1}{z},\quad
g_p(0)=p^{-1}i\log p,\\
\mathcal D_0(T)
=\sum_{\mu_{ij}\ne0}\sum_{p\in\{2,3\}}
 \frac{\mu_{ij}\overline{\xi_{p,ij}}}
 {|\xi_{2,ij}|^2+|\xi_{3,ij}|^2}
                     P_i\mathcal D_p(T)P_j.
\end{gathered}
\tag{MRR11}
\]
For the first equality, conjugation by $e^{zM}$ equals
$e^{z\operatorname{ad}_M}$: differentiating the conjugated
matrix yields the linear differential equation with that
generator, and its exponential series is its solution.
Since $\operatorname{ad}_M=\mathcal D_0+iI$, substituting
$z=i\log p$ proves the entire relation, including its
removable value at zero.
For the second equality, MRR3 shows that every denominator
in the stated range is strictly positive. On a given matrix
unit the two summands add to $\mu_{ij}$, whereas the common
kernel contributes zero. This proves the identity and makes
the relation of the finite-prime and infinitesimal defects
explicit in both directions.

\subsection{The actual nilpotent never belongs to this compatibility space}

The exact ST3 decomposition gives
$[M,R]=[C,R]$, because $[R,R]=0$.
The operator $\operatorname{ad}_C$ is self-adjoint on the
original Hilbert--Schmidt space. In fact cyclicity of the
finite trace gives
\[
\langle X,[C,Y]\rangle_{\rm HS,N}
=\operatorname{tr}(X^\dagger CY-X^\dagger YC)
=\langle[C,X],Y\rangle_{\rm HS,N}.
\tag{MRR12}
\]
The last equality uses $C=C^\dagger$ from ST3; it would not
follow by replacing $C$ with an unspecified amplitude shadow.
Consequently $\langle[C,R],R\rangle_{\rm HS,N}$ is real.
Expanding the squared norm, its cross term with $-iR$ has
zero real part. Since the rank-one norm is
$\|R\|_{\rm HS,N}^2=\epsilon^2\|f\|_N^2\|e\|_N^2=\epsilon^2$,
the exact obstruction is
\[
\boxed{
\|\mathcal D_0(R)\|_{\rm HS,N}^2
=\|[C,R]\|_{\rm HS,N}^2+\epsilon^2
\ge\epsilon^2>0.}
\tag{MRR13}
\]
In particular $R\notin\mathscr N$, including every resonance
in MRR4. The support nilpotent has not been discarded; its
nonzero image in the exact defect quotient MRR9 is exhibited.
By MRR11, its two-prime defect cannot vanish either.

There is a more explicit receiver in the unchanged original
rank-two plane and its orthogonal complement. Put
\[
\begin{gathered}
\Pi=ee^\dagger+ff^\dagger,\quad
\alpha=e^\dagger Ce\in\mathbb R,\quad
\beta=f^\dagger Cf\in\mathbb R,\quad c=e^\dagger Cf,\\
u=(I-\Pi)Ce,\quad v=(I-\Pi)Cf,\quad
Ce=\alpha e+\overline c f+u,\quad
Cf=c e+\beta f+v.
\end{gathered}
\tag{MRR14}
\]
Self-adjointness proves the conjugate coefficient in $Ce$.
Both $u$ and $v$ are orthogonal to $e$ and $f$.
Direct multiplication gives the entire original defect
\[
\boxed{
\mathcal D_0(R)
=\epsilon\bigl[c(ee^\dagger-ff^\dagger)
 +(\beta-\alpha-i)fe^\dagger+ve^\dagger-fu^\dagger\bigr].}
\tag{MRR15}
\]
For example $CR=\epsilon(Cf)e^\dagger$ and
$RC=\epsilon f(Ce)^\dagger$; substitution of MRR14 proves
every coefficient and sign.
The displayed four operator terms are mutually orthogonal
in the Hilbert--Schmidt metric. This follows from the
rank-one formula
$\langle xy^\dagger,zw^\dagger\rangle_{\rm HS,N}
=\langle x,z\rangle_N\langle w,y\rangle_N$
and the orthogonalities just stated; the diagonal difference
has squared norm two. Hence
\[
\boxed{
\|\mathcal D_0(R)\|_{\rm HS,N}^2
=\epsilon^2\bigl[1+(\beta-\alpha)^2+2|c|^2
                         +\|u\|_N^2+\|v\|_N^2\bigr].}
\tag{MRR16}
\]
It is also
$\epsilon^2[1+(\beta-\alpha)^2+
(\|Ce\|_N^2-\alpha^2)+(\|Cf\|_N^2-\beta^2)]$,
by MRR14. No source variance or cutoff term has been omitted.
Applying the actual defect to the original unit vector $e$
gives a stronger operator-norm bound:
\[
\begin{split}
\mathcal D_0(R)e&=\epsilon
                 \bigl(c e+(\beta-\alpha-i)f+v\bigr),\\
\|\mathcal D_0(R)\|_N
&\ge\epsilon\sqrt{|c|^2+(\beta-\alpha)^2+1+\|v\|_N^2}
\ge\epsilon.
\end{split}
\tag{MRR17}
\]
The three vectors in the first line are orthogonal, proving
the exact norm on that vector before taking the bound.

For completeness the metric distance of this actual $R$
from every compatible replacement is explicit and positive:
\[
\begin{gathered}
\operatorname{dist}_{\rm HS,N}(R,\mathscr N)^2
=\epsilon^2-b(R)^*\Gamma^{-1}b(R),\\
\operatorname{dist}_{\rm HS,N}(R,\mathscr N)
\ge\frac{\|\mathcal D_0(R)\|_{\rm HS,N}}{2\|M\|_N+1}
\ge\frac{\epsilon}{2\|M\|_N+1}>0.
\end{gathered}
\tag{MRR18}
\]
Use the off-resonance convention of MRR8 when the compatible
basis is empty. For any $T\in\mathscr N$,
$\mathcal D_0(R-T)=\mathcal D_0(R)$.
The original induced norms satisfy
$\|MX\|_{\rm HS,N},\|XM\|_{\rm HS,N}
\le\|M\|_N\|X\|_{\rm HS,N}$, proved by summing norms on
an original-metric orthonormal basis or by adjoints.
Thus $\|\mathcal D_0(X)\|_{\rm HS,N}
\le(2\|M\|_N+1)\|X\|_{\rm HS,N}$.
Taking the infimum proves the bound, with MRR13 giving
strict positivity.

\subsection{Complete observation blocks and all support labels}

Retain the original isometric section $J:B\to E$ from ST2,
$K=\ker\Lambda$, its isometric inclusion $\iota_K$, and
$W_0:B\oplus K\to E$, $W_0(b,z)=Jb+\iota_Kz$.
The original orthogonal decomposition proves that $W_0$ is
unitary for these unchanged metrics.
For $D=\mathcal D_0(R)$ the exact complete observation is
\[
W_0^\dagger D W_0=
\begin{pmatrix}
J^\dagger DJ&J^\dagger D\iota_K\\
\iota_K^\dagger DJ&\iota_K^\dagger D\iota_K
\end{pmatrix},\qquad
\|D\|_{\rm HS,N}^2
=\sum_{a,b\in\{B,K\}}\|D_{ab}\|_{\rm HS}^2.
\tag{MRR19}
\]
This follows by summing squared matrix entries in the
orthogonal source and target decomposition, and the inverse
map is $D=W_0(D_{ab})W_0^\dagger$.
Thus the actual nonzero obstruction MRR16 is retained by
these four blocks even if its visible $B,B$ compression
alone vanishes. The same statements hold for each component
of $\mathcal D_{2,3}(R)$ using MRR11.

For the full bounded distributive lattice $L$, give the
complex vector spaces $\mathcal E$, $\mathscr N$, their
quotient and each defect image their synchronized carriers
$V_L^\uparrow=\{(0,\lambda):\lambda\in L\}
\cup\{(v,1_L):v\in V\}$.
Every linear map above lifts by
$T^\uparrow(v,\lambda)=(Tv,\lambda)$.
Componentwise addition and support meet prove linearity
over $G_L(\mathbb C)$ and preservation of compositions.
In particular
\[
\begin{gathered}
(\mathcal D_0^\uparrow)^{-1}(Z_L)
=(\mathcal D_{2,3}^\uparrow)^{-1}(Z_L)
=\mathscr N_L^\uparrow,\\
(\mathcal E/\mathscr N)_L^\uparrow
\xrightarrow{\ \simeq\ }
(\operatorname{im}\mathcal D_{2,3})_L^\uparrow,
\qquad([T],\lambda)\longmapsto
                        (\mathcal D_{2,3}T,\lambda).
\end{gathered}
\tag{MRR20}
\]
In the first line $Z_L$ means the zero-amplitude fibre of
the appropriate synchronized target, including the direct
sum target for the pair. MRR3 and MRR9 prove the amplitude
statements; keeping $\lambda$ proves the displayed lifts.
For every $\lambda<1_L$, the quotient fibre over
$(0,\lambda)$ is its singleton zero state; over $(0,1_L)$
it is $\mathscr N\times\{1_L\}$.
The actual $R\ne0$ starts at top support and has nonzero
defect by MRR13, so its defect remains at that same top
support. All lower zero labels remain distinct and fixed
by the zero-amplitude action.

The original free-coordinate correspondence is the full
fibre from OLR17: in any retained basis of the positive
dimension $\dim\mathcal E=q^2$, synchronization sends
$(T,\boldsymbol\mu)$ to $(T,\bigvee_i\mu_i)$ with
$T_i\ne0\Rightarrow\mu_i=1_L$.
Its inverse fibre retains exactly those masks with that
join and those forced top coordinates.
Over the compatibility inverse image in MRR20 the amplitude
is exactly a member of MRR4 and all its permitted original
masks survive. No change of free-semimodule basis to the
spectral units has been asserted.
The metric pairing on every synchronized space is
$(\langle x,y\rangle,\lambda\wedge\mu)$, so MRR7--MRR19
retain their original amplitudes and every support meet.

The calculation therefore identifies the full finite space
compatible with Jannsen's stated geometric-Frobenius equation,
proves its exact relation to the two original prime operators,
and exhibits the complete defect of the actual programme
nilpotent. It does not replace that nilpotent by a resonant
matrix unit, identify its slot grades with arithmetic weights,
or infer a purity statement or RH from finite compatibility.

\subsection{Every resonance chain, break, and Jordan block}

Fix a resonance $\delta=1/(2m)$ with $1\le m\le k$, and
retain an arbitrary member of MRR4, with its full complex coefficients:
\[
\mathcal N=\sum_{a=0}^{k-m}\sum_{b=0}^k n_{ab}U_{ab,a+m,b}.
\]
For each $b\in\{0,\ldots,k\}$ and residue
$r\in\{0,\ldots,m-1\}$ define
\[
\begin{gathered}
L_r=\left\lfloor\frac{k-r}{m}\right\rfloor,\qquad
v_{r,b,j}=\ell_{r+jm,b}\quad(0\le j\le L_r),\\
c_{r,b,j}=n_{r+(j-1)m,b}\quad(1\le j\le L_r),\qquad
\mathcal Nv_{r,b,j}=c_{r,b,j}v_{r,b,j-1}\quad(j\ge1),\qquad
\mathcal Nv_{r,b,0}=0.
\end{gathered}
\tag{MRR21}
\]
Every original label occurs in exactly one such chain by
division with remainder of $a$ by $m$.
Thus this is the entire original object, with $b$ unchanged;
no two spectral labels are identified.
For the boundary tests below put
$c_{r,b,0}=c_{r,b,L_r+1}=0$. These two values mark the
absence of an edge outside the original chain; they are
not additional resonance coefficients.

A segment $[s,u]\subseteq[0,L_r]$ is one Jordan block
precisely when
\[
c_{r,b,s}=0,\qquad c_{r,b,u+1}=0,\qquad
c_{r,b,j}\ne0\quad(s<j\le u).
\tag{MRR22}
\]
Indeed a zero coefficient deletes exactly its one edge.
The remaining connected components are exactly these
segments, since the underlying finite graph is a path.
Their vector spans are $\mathcal N$-invariant and form
a direct sum, because they partition the original basis.
In one segment put
\[
w_{r,b,[s,u],h}=\mathcal N^h v_{r,b,u}
=\left(\prod_{j=u-h+1}^{u}c_{r,b,j}\right)v_{r,b,u-h}
\quad(0\le h\le u-s).
\tag{MRR23}
\]
The empty product is one. Every displayed product is nonzero
by MRR22, so these vectors are a basis of the segment.
They satisfy $\mathcal Nw_h=w_{h+1}$ up to its last vector,
which is sent to zero. Therefore the segment contributes
exactly one Jordan block of length $u-s+1$.
This proves the classification for every coefficient choice,
including all zero coefficients and all nonzero phases.
The vectors in MRR23 are actual powers of the original
operator; they have not been divided by their lengths or norms.

In particular the multiplicity of a Jordan block of length
$\ell\ge1$ is the completely specified integer
\[
\boxed{
\begin{aligned}
J_\ell(\mathcal N)
=\sum_{b=0}^k\sum_{r=0}^{m-1}
\sum_{s=0}^{L_r+1-\ell}
 &\mathbf1_{c_{r,b,s}=0}
  \mathbf1_{c_{r,b,s+\ell}=0}\\[-2pt]
 &\hspace{-20pt}\cdot
 \prod_{j=s+1}^{s+\ell-1}\mathbf1_{c_{r,b,j}\ne0}.
\end{aligned}}
\tag{MRR24}
\]
A sum with negative upper bound is empty.
Its term is one exactly for the segment in MRR22 whose
upper endpoint is $s+\ell-1$, and zero otherwise, proving
the formula. Every admissible Jordan type is thus exactly
the multiset obtained by partitioning each of the retained
chain lengths $L_r+1$ into ordered positive segment lengths,
independently for every $b$. Conversely any such partitions
are obtained by setting the separating coefficients to zero
and the remaining coefficients to any chosen nonzero complex
numbers, so this description is also sufficient.

For every integer $a\ge0$ the exact ranks and kernels are
\[
\begin{gathered}
\operatorname{rank}\mathcal N^a
=\sum_{\ell\ge1}J_\ell(\mathcal N)\max(\ell-a,0),\qquad
\dim\ker\mathcal N^a
=\sum_{\ell\ge1}J_\ell(\mathcal N)\min(a,\ell),\\
\dim\ker\mathcal N
=q-\#\{(a,b):n_{ab}\ne0\}.
\end{gathered}
\tag{MRR25}
\]
The first two identities follow by shifting each basis in
MRR23. The last follows because every nonzero coefficient
is one retained edge in a disjoint union of paths, whose
number of components is the number of vertices minus edges.
Equivalently the nonzero columns of $\mathcal N$ have
different target basis vectors and are independent.
The nilpotence index is the largest segment length, at most
$\lfloor k/m\rfloor+1$; the zero operator has index one.

The original spectral metric remains exact throughout.
Write $H_{ij}=\langle\ell_i,\ell_j\rangle_N$ as in MRR7.
For segments $A,B$, let $d_{A,h}$ be the literal product
in MRR23 and $i(A,h)$ its original remaining spectral label.
Then their entire Gram matrix, including different segments,
is
\[
\boxed{
\langle w_{A,h},w_{B,l}\rangle_N
=\overline{d_{A,h}}d_{B,l}H_{i(A,h),i(B,l)}.}
\tag{MRR26}
\]
This is substitution of MRR23 into the original sesquilinear
pairing. In particular distinct Jordan blocks are not assumed
orthogonal. The map from all these power vectors to the
original cardinal basis is a permutation followed by the
displayed nonzero diagonal coefficients; its inverse divides
each coefficient by that same product. This proves an exact
coordinate isomorphism and its inverse, with the metric given
by MRR26 rather than replaced by an identity matrix.

\subsection{The exact semidirect group carried by these chains}

For every positive real number $x$ retain the original
one-parameter extension
\[
F_x=\exp(i(\log x)M),\qquad
F_x\ell_{ab}=
x^{(2a-k)\delta}e^{i(2b-k)\gamma\log x}\ell_{ab},
\qquad F_xF_y=F_{xy}.
\tag{MRR27}
\]
The logarithm here is the unique real logarithm.
Its additive law and the commuting exponential series prove
the last equality. The original prime operators are exactly
the values at $x=p$.
For $\mathcal N$ in MRR4, the identical matrix-unit
calculation MRR2 gives
$F_x\mathcal NF_x^{-1}=x^{-1}\mathcal N$ for every $x>0$.

Define the set $\mathcal G=\mathbb R_{>0}\times\mathbb C$
with the following full operations:
\[
(x,t)(y,s)=(xy,t+x^{-1}s),\qquad
1_{\mathcal G}=(1,0),\qquad
(x,t)^{-1}=(x^{-1},-xt).
\tag{MRR28}
\]
Both associations of a triple have second coordinate
$t+x^{-1}s+(xy)^{-1}u$, proving associativity.
Substitution proves both the identity and the two inverse
identities. This is a topological group for the product
of the usual real and complex topologies, since these
displayed operations are continuous.
The exact representation is
\[
\boxed{
\rho_{\mathcal N}(x,t)=\exp(t\mathcal N)F_x,
\qquad
\rho_{\mathcal N}((x,t)(y,s))
=\rho_{\mathcal N}(x,t)\rho_{\mathcal N}(y,s).}
\tag{MRR29}
\]
Indeed conjugating the finite exponential by $F_x$ gives
$F_xe^{s\mathcal N}F_x^{-1}=e^{x^{-1}s\mathcal N}$,
term by term. Its product with $e^{t\mathcal N}$ is
$e^{(t+x^{-1}s)\mathcal N}$, since both exponents are scalar
multiples of the same operator. MRR27 proves the group law.
The identity maps to $I_E$, and the inverse maps to
$(e^{t\mathcal N}F_x)^{-1}$ by that law, with no change in
the order of the two original factors. Continuity follows
from the finite exponential polynomial and MRR27.

There is also a precise converse: for any endomorphism $T$,
the expression $e^{tT}F_x$ is a representation of MRR28
if and only if $F_xTF_x^{-1}=x^{-1}T$ for all $x>0$.
Necessity follows by evaluating the representation identity
on $(x,0)(1,s)=(x,x^{-1}s)$ and differentiating at $s=0$.
Sufficiency is the same exponential multiplication just
proved, valid for the convergent exponential even without
an a priori nilpotence assumption. Differentiating
$F_{e^u}TF_{e^u}^{-1}=e^{-u}T$ at $u=0$ gives
$i[M,T]=-T$, equivalently $[M,T]=iT$.
Conversely that commutator equation gives the conjugation
identity by exponentiation of $\operatorname{ad}_M$.
Together with MRR3--MRR4, this proves the full equivalence
\[
\boxed{
\rho_T\text{ is a representation of }\mathcal G
\ \Longleftrightarrow\
F_xTF_x^{-1}=x^{-1}T\ (x>0)
\ \Longleftrightarrow\ T\in\mathscr N.}
\tag{MRR30}
\]
For this original $M$, its last condition implies exactly
the nilpotent chain classification already proved.
Thus the representation has been constructed on the entire
object defined by the obstruction, with a converse
characterization rather than a merely formal resemblance.

The representation is faithful exactly when $\mathcal N\ne0$.
In the original cardinal ordering, $e^{t\mathcal N}$ has
all diagonal entries one and only strictly lowering
off-diagonal entries. Thus the diagonal of
$\rho_{\mathcal N}(x,t)$ is the original spectrum of $F_x$.
If it is the identity, any label with $2a-k\ne0$ gives
$x^{(2a-k)\delta}=1$, hence $x=1$ since $\delta>0$.
If $\mathcal N\ne0$ and $t\ne0$, the factorization
$e^{t\mathcal N}-I=t\mathcal N
\sum_{j\ge0}(t\mathcal N)^j/(j+1)!$
has invertible second factor, since it is $I$ plus a
nilpotent. Its value cannot be zero. Therefore $t=0$.
When $\mathcal N=0$ its exact kernel is
$\{(1,t):t\in\mathbb C\}$, by the same diagonal argument.

\subsection{Every phase, chain entry, and metric under the action}

For an original chain label $v_{r,b,j}$ put
$\omega_{r,b,j}(x)=
x^{(2(r+jm)-k)\delta}e^{i(2b-k)\gamma\log x}$.
The full action in the unchanged cardinal basis is
\[
\rho_{\mathcal N}(x,t)v_{r,b,j}
=\omega_{r,b,j}(x)
\sum_{h=0}^j\frac{t^h}{h!}
 \left(\prod_{l=j-h+1}^jc_{r,b,l}\right)v_{r,b,j-h}.
\tag{MRR31}
\]
This is the finite exponential applied after the original
$F_x$. Zero coefficients automatically terminate every
path at its actual break. The formula therefore works
without discarding any broken or length-one chain.

On one segment with top label $v_{r,b,u}$, let
$\omega_A(x)=\omega_{r,b,u}(x)$ and $d_A=u-s$.
The power basis in MRR23 has the exact diagonal action
$F_xw_{A,h}=\omega_A(x)x^{-h}w_{A,h}$, either by MRR27
or its conjugation law. Consequently
\[
\rho_{\mathcal N}(x,t)w_{A,h}
=\omega_A(x)x^{-h}
\sum_{l=h}^{d_A}\frac{t^{l-h}}{(l-h)!}w_{A,l}.
\tag{MRR32}
\]
This retains the segment's top spectral phase and every
power of $x$; its coefficients are not arithmetic weight
assignments. Combining MRR32 with the complete Gram
matrix $G_{A,l;B,n}$ from MRR26 gives the exact pairing
\[
\begin{split}
\langle\rho_{\mathcal N}(x,t)w_{A,h},
        \rho_{\mathcal N}(y,s)w_{B,j}\rangle_N
={}&\overline{\omega_A(x)}\omega_B(y)x^{-h}y^{-j}\\
&\cdot\sum_{l=h}^{d_A}\sum_{n=j}^{d_B}
\frac{\overline t^{\,l-h}s^{n-j}}
 {(l-h)!(n-j)!}\,G_{A,l;B,n}.
\end{split}
\tag{MRR33}
\]
This is finite sesquilinear expansion; it preserves cross
terms between distinct Jordan segments and all original
cutoff-dependent entries of $G_N$.

The uncentred mass has a full representation as well:
\[
\widetilde\rho_{\mathcal N}(x,t)
=x^{k/2}\rho_{\mathcal N}(x,t),\qquad
\det\rho_{\mathcal N}(x,t)=1,\qquad
\det\widetilde\rho_{\mathcal N}(x,t)=x^{kq/2}.
\tag{MRR34}
\]
The scalar $x^{k/2}$ is a group character of MRR28, so
its product with MRR29 obeys the same representation law.
Its values at $(p,0)$ are the actual $\widetilde F_p$.
The nilpotent exponential has determinant one by its
triangular diagonal. Also
$\sum_{a,b}\lambda_{ab}=0$ by separately summing
$2a-k$ and $2b-k$ from zero to $k$, so the product of
all eigenvalues of $F_x$ is one. This proves both
determinants, with the full original dimension $q$.

\subsection{Units, labelled zeros, and the full support correspondence}

Every operator in MRR29 is invertible, hence has the
exact full-lattice automorphism
\[
\rho_{\mathcal N,L}(x,t)(v,\lambda)
=(\rho_{\mathcal N}(x,t)v,\lambda),\qquad
\rho_{\mathcal N,L}(1,0)=I_{E_L^\uparrow}.
\tag{MRR35}
\]
Composition of lifts is composition of their amplitudes
with support unchanged. Thus MRR29 proves the group law,
and the lift at $(x^{-1},-xt)$ is its inverse, on every
original nonzero vector and on every $(0,\lambda)$.
Each labelled zero is fixed, individually. In particular
the identity fixes supported zero and actual absence as
their distinct states; it does not send them both to one
chosen zero.
Likewise $(\mathcal N_L^\uparrow)^a=(\mathcal N^a)_L^\uparrow$.
After the largest segment length, this is the supported-zero
operator $(v,\lambda)\mapsto(0,\lambda)$ rather than the
constant map at $\tau$. The exponential lift in MRR35
retains its identity term and remains an automorphism;
nilpotence of the amplitude does not make its exponential
a zero-valued character.
The pairing of any two lifted vectors is
$(\langle v,w\rangle_N,\lambda\wedge\mu)$, so the full
metric identities MRR26 and MRR33 carry exactly the same
support meet. Isometry is not assumed when the displayed
metric amplitudes change.

The correspondence with all original free masks can also
be stated with its exact identity. Let $\mathcal F_L$ be
the free-coordinate carrier in the original positive
dimension $q$, and let $P_{\rm syn}:\mathcal F_L\to E_L^\uparrow$
be the synchronization map of OLR17.
For $g\in\mathcal G$ define a relation
\[
\mathcal R_g=
\{(u,v)\in\mathcal F_L\times\mathcal F_L:
             P_{\rm syn}v=\rho_{\mathcal N,L}(g)P_{\rm syn}u\}.
\tag{MRR36}
\]
In original coordinates this means that the output
amplitude is exactly $\rho_{\mathcal N}(g)$ times the
input amplitude, their support joins agree, and every
nonzero coordinate forces its own support to $1_L$.
Every intermediate synchronized state has a free lift:
use all support coordinates equal to its label when its
amplitude is zero, and all equal to $1_L$ otherwise.
It follows, with relational composition in the indicated
order, that
\[
\mathcal R_g\circ\mathcal R_h=\mathcal R_{gh},\qquad
\mathcal R_g^{-1}=\mathcal R_{g^{-1}},\qquad
\mathcal R_{1_{\mathcal G}}
=\{(u,v):P_{\rm syn}u=P_{\rm syn}v\}.
\tag{MRR37}
\]
For the forward containment in the first equality,
apply the representation law to any intermediate state.
For the reverse containment, choose a free lift of
$\rho_{\mathcal N,L}(h)P_{\rm syn}u$ as the intermediate
state; the group law puts the prescribed final state
in the second relation. Inverting the defining equality
proves the second identity, and MRR35 proves the third.
Thus the identity relation on free masks is precisely
equality after synchronization, not the diagonal relation
on individual masks. The quotient by that equivalence
is exactly the synchronized carrier, on which MRR35 is
the proved group action with its genuine identity.
This gives the full free-mask correspondence without
claiming a new free-semimodule basis change or deleting
any original mask.

MRR21--MRR37 construct the chains, group, representation,
metric and full-support correspondences determined by
the resonance space. The actual ST3 boundary nilpotent
$R$ remains outside that space by MRR13--MRR17; none of
these constructions silently replaces it by a chosen
compatible coefficient pattern or supplies an arithmetic
purity identification.
